\documentclass[11pt]{article}
\usepackage[T1]{fontenc}
\usepackage[margin=1in]{geometry}
\usepackage{amsmath,amssymb,amsthm,mathtools}
\usepackage{graphicx}
\usepackage{microtype}
\usepackage[hidelinks]{hyperref}

\newtheorem{theorem}{Theorem}
\newtheorem{lemma}[theorem]{Lemma}
\newtheorem{remark}[theorem]{Remark}
\newcommand{\g}{\mathfrak g}
\newcommand{\h}{\mathfrak h}
\newcommand{\R}{\mathbb R}

\title{Stabilizers of smooth Banach--Lie actions\\are Lie subgroups}
\author{Autonomous solution packet for arXiv:math/0703710\\Model: GPT5.6}
\date{11 August 2026}

\begin{document}
\maketitle

\begin{abstract}
An--Neeb ask whether the stabilizer of a point under an arbitrary smooth
action of a Banach--Lie group on a Banach manifold is always a Banach--Lie
subgroup.  We prove that it is; in fact a $C^2$ action suffices.  The key is a
uniform no-small-period lemma.  In a chart at the fixed point, the vector
field generated by $x\in\g$ has local Lipschitz constant $O(\|x\|)$.  If
$\exp(x)$ fixes the point and $x$ is small, Gronwall's inequality forces the
initial velocity to vanish.  Flow uniqueness then makes the whole
one-parameter group fix the point.  Consequently the stabilizer is, in an
exponential chart, exactly its closed infinitesimal stabilizer algebra.
\end{abstract}

\section{Source problem and answer}

The source proves the conclusion when the tangent space to the orbit is
closed and then asks the unrestricted question \cite[Problem 4.5]{AN}:

\begin{center}
\includegraphics[width=.96\textwidth]{figures/open_problem_crop.png}
\end{center}

Here ``Lie subgroup'' means a Banach--Lie group in the subspace topology,
without requiring its Lie algebra to be complemented in $\g$.

\begin{theorem}[affirmative answer]\label{thm:main}
Let $G$ be a Banach--Lie group acting by a $C^2$ action
\[
 \sigma:G\times M\longrightarrow M
\]
on a Banach $C^2$-manifold $M$.  For every $p\in M$, the stabilizer
\[
 G_p=\{g\in G:g.p=p\}
\]
is a Banach--Lie subgroup of $G$ with its subspace topology.
\end{theorem}

Thus Problem 4.5 has a positive answer, with weaker regularity than the
smoothness assumed there.  No closed-range or complementability hypothesis is
needed.

\section{Proof intuition}

Write $x_M$ for the fundamental vector field generated by $x\in\g$.  Near
$p$, its chart expression $F(x,y)$ is linear in $x$.  Its derivative in the
$y$-variable is consequently bounded by $C\|x\|$, uniformly near $p$.
The orbit $t\mapsto\exp(tx).p$ solves $y'=F(x,y)$.  If it returns to $p$ at
time one, comparison with the constant vector $F(x,0)$ shows that
\[
 \|F(x,0)\|\leq C\|x\|e^{C\|x\|}\|F(x,0)\|.
\]
For small $x$ the right coefficient is less than one, so the initial velocity
vanishes.  The orbit is then constant by ODE uniqueness.  This converts a
time-one stabilizer relation into an infinitesimal stabilizer relation,
uniformly in every sufficiently small logarithm.

\section{A parametric no-small-return lemma}

\begin{lemma}\label{lem:return}
Let $P$ and $E$ be Banach spaces, let $V\subseteq E$ be a neighborhood of
$0$, let $A:V\to\mathcal L(P,E)$ be $C^1$, and put
$F(x,y)=A(y)x$.
For $x$ near $0$, suppose $y_x:[0,1]\to V$ solves
\[
 y_x'(t)=F(x,y_x(t)),\qquad y_x(0)=0,
\]
and that all these curves remain in one fixed neighborhood of $0$ in $V$.
Then there is a neighborhood $U$ of $0$ in $P$ such that
\[
 x\in U,\quad y_x(1)=0 \quad\Longrightarrow\quad F(x,0)=0.
\]
In that event $y_x(t)=0$ for every $t\in[0,1]$.
\end{lemma}

\begin{proof}
The derivative of $A$ gives a continuous family of bounded bilinear maps
\[
 B(y)(x,z):=(DA(y)z)x.
\]
After shrinking the fixed $y$-neighborhood, there is $C>0$ such that
$\|B(y)\|\leq C$ throughout it.  The mean-value formula on a convex ball
therefore gives
\begin{equation}\label{eq:lipschitz}
 \|F(x,y)-F(x,0)\|\leq C\|x\|\,\|y\|.
\end{equation}

Put $v=F(x,0)$ and $\lambda=C\|x\|$.  Integrating the ODE and using
\eqref{eq:lipschitz} gives, for $0\leq t\leq1$,
\[
 \|y_x(t)\|\leq t\|v\|+\lambda\int_0^t\|y_x(s)\|\,ds.
\]
Gronwall's inequality yields
\begin{equation}\label{eq:gronwall}
 \|y_x(t)\|\leq t e^{\lambda t}\|v\|.
\end{equation}
If $y_x(1)=0$, another integration of the ODE gives
\[
 v=-\int_0^1\big(F(x,y_x(s))-F(x,0)\big)\,ds.
\]
Combining this with \eqref{eq:lipschitz} and \eqref{eq:gronwall},
\[
 \|v\|
 \leq\lambda\int_0^1s e^{\lambda s}\,ds\,\|v\|
 \leq\lambda e^\lambda\|v\|.
\]
Shrink $U$ so that $\lambda e^\lambda<1$ there.  Then $v=0$.  The constant
curve is now a solution with the same initial value, so uniqueness for Banach
space ODEs gives $y_x\equiv0$.
\end{proof}

\section{Proof of the theorem}

\begin{proof}[Proof of Theorem~\ref{thm:main}]
Fix $p\in M$ and put $H=G_p$.  It is a closed subgroup.  Let
\[
 \alpha_p:\g\longrightarrow T_pM,
 \qquad \alpha_p(x)=\left.\frac d{dt}\right|_{t=0}\exp_G(tx).p,
\]
and set $\h=\ker\alpha_p$.  This is a closed linear subspace.  It is also a
Lie subalgebra: the fundamental vector fields generated by elements of
$\h$ vanish at $p$, and hence so does their bracket.

Choose a chart $\kappa:W\to V\subseteq E$ with $\kappa(p)=0$.  The chart
expression of the fundamental fields is
\[
 F(x,y)=T_{\kappa^{-1}(y)}\kappa
 \left(\left.\frac d{dt}\right|_{t=0}
       \exp_G(tx).\kappa^{-1}(y)\right).
\]
It has the form $F(x,y)=A(y)x$ for a $C^1$ map
$A:V\to\mathcal L(\g,E)$: the $C^2$ regularity of the action supplies
precisely this $C^1$ dependence of its first differential.  Thus
Lemma~\ref{lem:return} applies.

For all sufficiently small $x\in\g$, the curve
\[
 y_x(t)=\kappa(\exp_G(tx).p),\qquad 0\leq t\leq1,
\]
stays in one fixed chart neighborhood.  This follows uniformly in $t$ from
continuity of $(t,x)\mapsto\exp_G(tx).p$ and compactness of $[0,1]$.
It solves the ODE of Lemma~\ref{lem:return}, with
$F(x,0)=T_p\kappa(\alpha_p(x))$.

Take a sufficiently small open $0$-neighborhood $U\subseteq\g$ on which
$\exp_G$ is a diffeomorphism onto an identity neighborhood and on which the
lemma applies.  If $x\in U$ and $\exp_G(x)\in H$, then $y_x(1)=0$, so the
lemma implies $\alpha_p(x)=0$, i.e. $x\in\h$.  Conversely, if $x\in\h$,
the orbit through $p$ and the constant curve through $p$ solve the same
autonomous ODE with the same initial value.  Uniqueness gives
$\exp_G(tx).p=p$ for all $t\in\R$, and in particular $\exp_G(x)\in H$.
We have proved the local identity
\begin{equation}\label{eq:local}
 H\cap\exp_G(U)=\exp_G(U\cap\h).
\end{equation}

Since $\h$ is closed, it is a Banach space.  Equation~\eqref{eq:local}
makes $\log_G$ a chart from the open identity neighborhood
$H\cap\exp_G(U)$ in the subspace topology onto $U\cap\h$.  Its left
translates form a Banach-manifold atlas on $H$.  The chart expressions of
multiplication and inversion are restrictions of the smooth local group
operations of $G$ (equivalently, of the local BCH operations), so they are
smooth.  Thus $H$ is a Banach--Lie group, and its inclusion into $G$ is a
smooth topological embedding.  Hence it is a Banach--Lie subgroup in the
source's sense.
\end{proof}

\begin{remark}[what is and is not solved]
The theorem removes the closed-orbit-tangent assumption in An--Neeb's
Theorem 4.2 and fully answers their Problem 4.5.  It does not construct a
Banach-manifold structure on $G/H$ for an arbitrary nonsplit Banach--Lie
subgroup $H$, so the opposite direction of their Conjecture 4.1 remains
outside the claim.
\end{remark}

\section{Literature and novelty boundary}

Jotz--Neeb prove the proper-action case by showing that every orbit tangent
has closed range \cite{JN}; combined with An--Neeb's theorem, this yields Lie
stabilizers for proper actions.  Diez--Rudolph still state that the general
stabilizer question is unknown even for Banach--Lie actions
\cite{DR}.  Most decisively, Gl\"ockner--Neeb's 2026 monograph manuscript
repeats the exact Banach transformation-group problem and the same known
closed-range and linear-action subcases \cite{GN}.

The run searched its cheap indexes and performed bounded arXiv searches on
2026-08-11 using the exact problem sentence and combinations of
``stabilizer'', ``Banach--Lie subgroup'', ``smooth action'', and
``Problem IX.3.b''.  These searches returned the source, the proper-action
paper, the 2022 restatement, and the 2026 monograph, but no unrestricted
answer.  The conservative status is therefore a candidate full solution
pending human verification.  Review should focus on the uniform mixed-
derivative estimate in Lemma~\ref{lem:return}; it is the step that rules out
the natural weighted-rotation counterexample attempts.

\begin{thebibliography}{9}
\bibitem{AN}
J. An and K.-H. Neeb,
\emph{An implicit function theorem for Banach spaces and some applications},
Math. Z. 262 (2009), 627--643; arXiv:math/0703710.
\url{https://arxiv.org/abs/math/0703710}

\bibitem{JN}
M. Jotz and K.-H. Neeb,
\emph{Closedness of the tangent spaces to the orbits of proper actions},
J. Lie Theory 18 (2008), 517--521; arXiv:0804.4858.
\url{https://arxiv.org/abs/0804.4858}

\bibitem{DR}
T. Diez and G. Rudolph,
\emph{Normal form of equivariant maps in infinite dimensions},
Ann. Global Anal. Geom. 61 (2022), 159--213; arXiv:2010.10165.
\url{https://arxiv.org/abs/2010.10165}

\bibitem{GN}
H. Gl\"ockner and K.-H. Neeb,
\emph{Infinite-Dimensional Lie Groups}, version of 2 January 2026;
arXiv:2602.12362.
\url{https://arxiv.org/abs/2602.12362}
\end{thebibliography}

\end{document}
